\documentclass{article}
\usepackage[utf8]{inputenc}
\title{Early rehabilitation after orthopaedic and cardiac surgery: what recent pilot trials show}
\author{}
\date{}
\begin{document}
\maketitle

\section{Introduction}
Pilot trials are the usual first step for rehabilitation interventions after surgery. After total joint replacement, all participants in the REPAIRS pilot received a non-opioid regimen of naproxen \SI{500}{\micro\gram} every 12 hours \cite{yang2025repairs}, and the trial recruited 53 participants, a recruitment rate of 41\% \cite{yang2025repairs}. In older adults after cardiac surgery, 100 of 336 patients screened were randomized to home exercise training or usual activity \cite{tamuleviciute2026}. Modified sit-to-stand training increased loading on the affected limb after hip fracture \cite{noda2022sts}, and home-based eccentric training after hip arthroplasty was feasible in adults over 75 \cite{lindgren2025}.

Technology-assisted approaches are spreading. In robot-assisted therapy for adhesive capsulitis, external rotation declined over time in the treatment group \cite{ryu2026robot}. Virtual-reality rehabilitation was tested for postoperative C5 palsy \cite{suk2025vr}, and technology-assisted upper-extremity training after spinal cord injury in a crossover pilot \cite{kilkki2025}.

\section{Perioperative context}
Hypotension prevention during hip fracture surgery \cite{luney2025}, outcomes of subtotal cholecystectomy techniques \cite{aloraini2025}, prosthetic joint infection after hemiarthroplasty \cite{mahmoud2025}, and upstaging after surgery for renal cell carcinoma \cite{goodstein2025} define the clinical background for these programmes.

\bibliographystyle{vancouver}
\bibliography{references}
\end{document}
